\documentclass[a4paper,10pt]{article}

\usepackage[utf8]{inputenc}
\usepackage[T1]{fontenc}
\usepackage[spanish]{babel}
\usepackage[margin=0.85in]{geometry}
\usepackage{titlesec}
\usepackage{enumitem}
\usepackage{hyperref}
\usepackage{xcolor}
\usepackage{indentfirst}

\titleformat{\section}{\color{blue}\large\bfseries}{}{0em}{}[\titlerule]
\pagenumbering{gobble}
\setlist[itemize]{leftmargin=1.5em, itemsep=0.15em, topsep=0.2em}

\begin{document}

\begin{center}
    {\LARGE \textbf{Pedro Henrique Klein}}\\
    \vspace{0.2cm}
    \href{mailto:pedrohenriquekleinphg@gmail.com}{pedrohenriquekleinphg@gmail.com} \textbar\
    (54) 99704-9008 \textbar\
    \href{https://www.linkedin.com/in/pedro-henrique-klein-a41122221/}{LinkedIn} \textbar\
    \href{https://github.com/Pedrinhonitz}{GitHub} \textbar\
    Chapecó, SC, Brasil
\end{center}

\vspace{0.3cm}

\section*{Sobre mí}

Empecé a programar a los 15 años, primero con Arduino y después con Python, por mi cuenta. A los 16 entré al mercado como desarrollador y, hasta los 18, trabajé en backend web con Django, Node.js, PHP, JavaScript y TypeScript.

En 2021 ingresé a Ciencias de la Computación en la Universidade Federal da Fronteira Sul (UFFS), en Chapecó, donde sigo estudiando. A inicios de 2022 dejé el backend y pasé a ingeniería de datos. Desde entonces construyo pipelines, modelo datos y orquestro procesos con Python, SQL, Airflow, GCP y AWS.

Esa mezcla de desarrollo y datos me permite transitar entre áreas técnicas y trabajar bien con equipos distintos.

\section*{Formación académica}

\textbf{Universidade Federal da Fronteira Sul (UFFS)} \hfill 2021 a 2027 \\
Licenciatura en Ciencias de la Computación, campus Chapecó

\section*{Experiencia profesional}

\textbf{BIP Brasil, proyecto Pro Carbono (Bayer)} \hfill Sep 2024 a Ago 2026

\textit{Ingeniero de Datos Senior} \hfill Ago 2025 a Ago 2026
\begin{itemize}
    \item Implemento y mantengo pipelines con Apache Airflow, PySpark, AWS Glue y Astronomer Cosmos.
    \item Creo aplicaciones y APIs internas con Python, FastAPI y TypeScript para automatizar procesos e integrar sistemas.
    \item Modelo datos con DBT, con calidad, estructura y trazabilidad de las transformaciones.
    \item Administro y optimizo entornos Linux, Docker y Kubernetes para cargas de ingeniería de datos.
    \item Escribo y optimizo consultas en PostgreSQL, Amazon Athena y BigQuery para grandes volúmenes de datos.
    \item Armo pipelines de CI/CD, reviso código y defino prácticas de ingeniería.
    \item Contribuyo a Apache Airflow con correcciones y mejoras para la comunidad.
\end{itemize}

\textit{Ingeniero de Datos Pleno} \hfill Sep 2024 a Ago 2025
\begin{itemize}
    \item Implementé y mantuve pipelines con Apache Airflow y PySpark.
    \item Construí soluciones de procesamiento y almacenamiento en AWS y Google Cloud Platform.
    \item Modelé datos con DBT y Elementary.
    \item Implanté CI/CD con GitHub Actions, Jenkins y Terraform, además de estándares de código y revisión de pull requests.
\end{itemize}

\textbf{Descomplica} \hfill Ene 2022 a Sep 2024 \\
\textit{Ingeniero de Datos Pleno}
\begin{itemize}
    \item Orquesto flujos en Apache Airflow para ejecutar y monitorear pipelines de datos.
    \item Empaqueto cargas con Docker y las orquestro en Kubernetes.
    \item Administro y optimizo entornos Linux.
    \item Construyo soluciones de data warehouse con BigQuery, Amazon Athena y Snowflake.
    \item Armo dashboards e informes en Metabase.
    \item Automatizo ETL con Python y Shell.
    \item Escribo y optimizo SQL para grandes volúmenes de datos.
\end{itemize}

\textbf{IXC-Soft} \hfill Nov 2021 a Ene 2022 \\
\textit{Pasante en Ingeniería de Software}
\begin{itemize}
    \item Desarrollé backend en PHP.
    \item Armé interfaces con Twig.
    \item Empaqueté entornos de desarrollo con Docker.
    \item Administré y optimicé bases MariaDB.
\end{itemize}

\section*{Open Source}

\textbf{\href{https://github.com/apache/airflow/pull/66431}{Apache Airflow (Apache Software Foundation)}} \\
Pull request aceptado en el repositorio oficial.
\begin{itemize}
    \item Corregí el scheduler en el tratamiento de \textit{deferred tasks}, eliminando una condición de carrera con eventos obsoletos enviados por el executor.
    \item Trabajé con los maintainers en la revisión técnica hasta que el cambio fue aprobado e integrado al código principal.
\end{itemize}

\section*{Habilidades técnicas}

\begin{itemize}
    \item \textbf{Lenguajes:} Python, SQL, TypeScript, JavaScript, Bash, Go, Java, C, C\texttt{++}, Haskell y Ruby.
    \item \textbf{Ingeniería de datos:} Apache Airflow, Astronomer Cosmos, DBT, Elementary, PySpark, AWS Glue, BigQuery, Snowflake, Amazon Athena, Luigi y Dagster.
    \item \textbf{DevOps e infraestructura:} Docker, Kubernetes, Terraform, GitHub Actions, Jenkins y Linux.
    \item \textbf{Bases de datos:} PostgreSQL, SQL Server, Oracle y Neo4j.
    \item \textbf{Cloud:} Amazon Web Services (AWS) y Google Cloud Platform (GCP).
\end{itemize}

\section*{Idiomas}

\begin{itemize}
    \item Portugués: nativo
    \item Inglés: A2
    \item Español: B1
\end{itemize}

\end{document}
